% Options for packages loaded elsewhere
% Options for packages loaded elsewhere
\PassOptionsToPackage{unicode}{hyperref}
\PassOptionsToPackage{hyphens}{url}
\PassOptionsToPackage{dvipsnames,svgnames,x11names}{xcolor}
%
\documentclass[
  british,
]{article}
\usepackage{xcolor}
\usepackage{amsmath,amssymb}
\setcounter{secnumdepth}{5}
\usepackage{iftex}
\ifPDFTeX
  \usepackage[T1]{fontenc}
  \usepackage[utf8]{inputenc}
  \usepackage{textcomp} % provide euro and other symbols
\else % if luatex or xetex
  \usepackage{unicode-math} % this also loads fontspec
  \defaultfontfeatures{Scale=MatchLowercase}
  \defaultfontfeatures[\rmfamily]{Ligatures=TeX,Scale=1}
\fi
\usepackage{lmodern}
\ifPDFTeX\else
  % xetex/luatex font selection
\fi
% Use upquote if available, for straight quotes in verbatim environments
\IfFileExists{upquote.sty}{\usepackage{upquote}}{}
\IfFileExists{microtype.sty}{% use microtype if available
  \usepackage[]{microtype}
  \UseMicrotypeSet[protrusion]{basicmath} % disable protrusion for tt fonts
}{}
\makeatletter
\@ifundefined{KOMAClassName}{% if non-KOMA class
  \IfFileExists{parskip.sty}{%
    \usepackage{parskip}
  }{% else
    \setlength{\parindent}{0pt}
    \setlength{\parskip}{6pt plus 2pt minus 1pt}}
}{% if KOMA class
  \KOMAoptions{parskip=half}}
\makeatother
% Make \paragraph and \subparagraph free-standing
\makeatletter
\ifx\paragraph\undefined\else
  \let\oldparagraph\paragraph
  \renewcommand{\paragraph}{
    \@ifstar
      \xxxParagraphStar
      \xxxParagraphNoStar
  }
  \newcommand{\xxxParagraphStar}[1]{\oldparagraph*{#1}\mbox{}}
  \newcommand{\xxxParagraphNoStar}[1]{\oldparagraph{#1}\mbox{}}
\fi
\ifx\subparagraph\undefined\else
  \let\oldsubparagraph\subparagraph
  \renewcommand{\subparagraph}{
    \@ifstar
      \xxxSubParagraphStar
      \xxxSubParagraphNoStar
  }
  \newcommand{\xxxSubParagraphStar}[1]{\oldsubparagraph*{#1}\mbox{}}
  \newcommand{\xxxSubParagraphNoStar}[1]{\oldsubparagraph{#1}\mbox{}}
\fi
\makeatother


\usepackage{longtable,booktabs,array}
\newcounter{none} % for unnumbered tables
\usepackage{calc} % for calculating minipage widths
% Correct order of tables after \paragraph or \subparagraph
\usepackage{etoolbox}
\makeatletter
\patchcmd\longtable{\par}{\if@noskipsec\mbox{}\fi\par}{}{}
\makeatother
% Allow footnotes in longtable head/foot
\IfFileExists{footnotehyper.sty}{\usepackage{footnotehyper}}{\usepackage{footnote}}
\makesavenoteenv{longtable}
\usepackage{graphicx}
\makeatletter
\newsavebox\pandoc@box
\newcommand*\pandocbounded[1]{% scales image to fit in text height/width
  \sbox\pandoc@box{#1}%
  \Gscale@div\@tempa{\textheight}{\dimexpr\ht\pandoc@box+\dp\pandoc@box\relax}%
  \Gscale@div\@tempb{\linewidth}{\wd\pandoc@box}%
  \ifdim\@tempb\p@<\@tempa\p@\let\@tempa\@tempb\fi% select the smaller of both
  \ifdim\@tempa\p@<\p@\scalebox{\@tempa}{\usebox\pandoc@box}%
  \else\usebox{\pandoc@box}%
  \fi%
}
% Set default figure placement to htbp
\def\fps@figure{htbp}
\makeatother


% definitions for citeproc citations
\NewDocumentCommand\citeproctext{}{}
\NewDocumentCommand\citeproc{mm}{%
  \begingroup\def\citeproctext{#2}\cite{#1}\endgroup}
\makeatletter
 % allow citations to break across lines
 \let\@cite@ofmt\@firstofone
 % avoid brackets around text for \cite:
 \def\@biblabel#1{}
 \def\@cite#1#2{{#1\if@tempswa , #2\fi}}
\makeatother
\newlength{\cslhangindent}
\setlength{\cslhangindent}{1.5em}
\newlength{\csllabelwidth}
\setlength{\csllabelwidth}{3em}
\newenvironment{CSLReferences}[2] % #1 hanging-indent, #2 entry-spacing
 {\begin{list}{}{%
  \setlength{\itemindent}{0pt}
  \setlength{\leftmargin}{0pt}
  \setlength{\parsep}{0pt}
  % turn on hanging indent if param 1 is 1
  \ifodd #1
   \setlength{\leftmargin}{\cslhangindent}
   \setlength{\itemindent}{-1\cslhangindent}
  \fi
  % set entry spacing
  \setlength{\itemsep}{#2\baselineskip}}}
 {\end{list}}
\usepackage{calc}
\newcommand{\CSLBlock}[1]{\hfill\break\parbox[t]{\linewidth}{\strut\ignorespaces#1\strut}}
\newcommand{\CSLLeftMargin}[1]{\parbox[t]{\csllabelwidth}{\strut#1\strut}}
\newcommand{\CSLRightInline}[1]{\parbox[t]{\linewidth - \csllabelwidth}{\strut#1\strut}}
\newcommand{\CSLIndent}[1]{\hspace{\cslhangindent}#1}

\ifLuaTeX
\usepackage[bidi=basic,shorthands=off]{babel}
\else
\usepackage[bidi=default,shorthands=off]{babel}
\fi
\ifLuaTeX
  \usepackage{selnolig} % disable illegal ligatures
\fi


\setlength{\emergencystretch}{3em} % prevent overfull lines

\providecommand{\tightlist}{%
  \setlength{\itemsep}{0pt}\setlength{\parskip}{0pt}}



 


% Get Math Fonts
\usepackage{fontspec}
\usepackage{unicode-math}
\setmathfont{Latin Modern Math}
\setmathfont[range={\mathscr,\mathbfscr,\mathbb}]{XITSMath-Regular}
[    Extension = .otf,
      BoldFont = XITSMath-Bold
]
\makeatletter
\@ifpackageloaded{caption}{}{\usepackage{caption}}
\AtBeginDocument{%
\ifdefined\contentsname
  \renewcommand*\contentsname{Table of contents}
\else
  \newcommand\contentsname{Table of contents}
\fi
\ifdefined\listfigurename
  \renewcommand*\listfigurename{List of Figures}
\else
  \newcommand\listfigurename{List of Figures}
\fi
\ifdefined\listtablename
  \renewcommand*\listtablename{List of Tables}
\else
  \newcommand\listtablename{List of Tables}
\fi
\ifdefined\figurename
  \renewcommand*\figurename{Figure}
\else
  \newcommand\figurename{Figure}
\fi
\ifdefined\tablename
  \renewcommand*\tablename{Table}
\else
  \newcommand\tablename{Table}
\fi
}
\@ifpackageloaded{float}{}{\usepackage{float}}
\floatstyle{ruled}
\@ifundefined{c@chapter}{\newfloat{codelisting}{h}{lop}}{\newfloat{codelisting}{h}{lop}[chapter]}
\floatname{codelisting}{Listing}
\newcommand*\listoflistings{\listof{codelisting}{List of Listings}}
\usepackage{amsthm}
\theoremstyle{definition}
\newtheorem{definition}{Definition}[section]
\theoremstyle{plain}
\newtheorem{proposition}{Proposition}[section]
\theoremstyle{plain}
\newtheorem{theorem}{Theorem}[section]
\theoremstyle{remark}
\AtBeginDocument{\renewcommand*{\proofname}{Proof}}
\newtheorem*{remark}{Remark}
\newtheorem*{solution}{Solution}
\newtheorem{refremark}{Remark}[section]
\newtheorem{refsolution}{Solution}[section]
\makeatother
\makeatletter
\makeatother
\makeatletter
\@ifpackageloaded{caption}{}{\usepackage{caption}}
\@ifpackageloaded{subcaption}{}{\usepackage{subcaption}}
\makeatother

\newcommand{\AmenableGroup}{{\Gamma}}
\newcommand{\AmenableGroupElement}{{\gamma}}
\newcommand{\ContinuousFunctionSpace}[1]{{\mathscr{C}({#1})}}
\newcommand{\CountingMeasure}[1]{{\left|#1\right|}}
\newcommand{\Folner}[1][\vphantom{}]{{\Phi}_{#1}}
\newcommand{\FurstenbergFolner}[1][{\vphantom{}}]{{\Psi}_{#1}}
\newcommand{\Group}{{\Gamma}}
\newcommand{\GroupAction}[2]{T^{#1}{#2}}
\newcommand{\KroneckerAction}[2]{S^{#1}{#2}}
\newcommand{\KroneckerFactorMap}{p}
\newcommand{\KroneckerGroup}{{H}}
\newcommand{\KroneckerMeasure}{{m_Z}}
\newcommand{\KroneckerProductAction}[2]{(T\times S)^{#1}{#2}}
\newcommand{\KroneckerSpace}{{Z}}
\newcommand{\KroneckerSpaceElement}{{z}}
\newcommand{\Measure}{{\mu}}
\newcommand{\ProjectionMap}[1][\vphantom{}]{{\pi}_{#1}}
\newcommand{\ShiftSpace}{{X}}
\newcommand{\ShiftSpaceElement}{{x}}
\newcommand{\Topology}{{\tau}}
\newcommand{\FurstenbergSystem}{(\ShiftSpace,\Measure,T)}
\newcommand{\KroneckerFactor}{{(\KroneckerSpace,\KroneckerMeasure,S)}}

\usepackage{bookmark}
\IfFileExists{xurl.sty}{\usepackage{xurl}}{} % add URL line breaks if available
\urlstyle{same}
% fallback for those not using the hyperref driver hyperxmp:
\makeatletter
\@ifundefined{xmpquote}{\newcommand{\xmpquote}[1]{#1}}{}
\makeatother
\hypersetup{
  pdftitle={Kronecker Factor},
  pdflang={en-GB},
  colorlinks=true,
  linkcolor={blue},
  filecolor={Maroon},
  citecolor={Blue},
  urlcolor={Blue},
  pdfcreator={LaTeX via pandoc}}


\title{Kronecker Factor}
\author{}
\date{2026-09-26}
\begin{document}
\maketitle


We want to prove the following result generalised for amenable groups,
where \(T\) is replaced with \(\AmenableGroup\):

\begin{proposition}[{Host, 2019, Proposition
5}]\protect\hypertarget{prp-KroneckerGenericSource0}{}\label{prp-KroneckerGenericSource0}

Let \((\ShiftSpace,T)\) be a topological dynamical system,
\(\ShiftSpaceElement_0\in\ShiftSpace\), and \(\Measure\) be an ergodic
invariant probability measure supported on the closed orbit of
\(\ShiftSpaceElement_0\) under \(T\). Let
\((\KroneckerFactor,\KroneckerMeasure,S)\) be the Kronecker factor of
\((\ShiftSpace,\Measure,T)\), with factor map
\(\KroneckerFactorMap:\ShiftSpace\rightarrow\KroneckerSpace\). Let
\(\ShiftSpace\times\KroneckerSpace\) be endowed with the transformation
of \(T\times S\). Let \(\tilde\Measure\) be the measure on
\(\ShiftSpace\times\KroneckerSpace\) and image of \(\Measure\) under the
map \(\ShiftSpace\rightarrow\ShiftSpace\times\KroneckerSpace\) where
\(\ShiftSpaceElement\mapsto(\ShiftSpaceElement,\KroneckerFactorMap(\ShiftSpaceElement))\).

Then there exists a Følner sequence \(\tilde{\Folner}\) and a point
\(\KroneckerSpaceElement_0\in\KroneckerSpace\) such that
\((\ShiftSpaceElement_0,\KroneckerSpaceElement_0)\) is generic for
\(\tilde{\Measure}\) along \(\tilde{\Folner}\).

\end{proposition}

Host (\citeproc{ref-host2019short}{2019}) also concludes that, for every
\(f\in\ContinuousFunctionSpace{X}\) and every
\(h\in\ContinuousFunctionSpace{Z}\), we have
\[\int_Xf(\ShiftSpaceElement)\cdot (h\circ\KroneckerFactorMap)(\ShiftSpaceElement)=\lim_{N\rightarrow\infty}\frac{1}{\CountingMeasure{\tilde{\Folner}_N}}\sum_{n\in\tilde{\Folner}_N}f(T^n\ShiftSpaceElement_0)h(S^n\KroneckerSpaceElement_0). \]

\section{Background}\label{background}

Firstly, we have the defininitions of the Kronecker subsystem and
Kronecker factor below:

\begin{definition}[{Jamneshan and Kreidler, 2025, Definition 7.1.14 \&
Subsection
13.3}]\protect\hypertarget{def-KroneckerFactor}{}\label{def-KroneckerFactor}

Let \((X,T)\) be a measure-preserving system. We then write
\(J_\text{kro}:(X_\text{kro},T_\text{kro})\rightarrow(X,T)\) for the
extension \(J_E:(X_E,T_E)\rightarrow(X,T)\) defined by the invariant
Markov sublattice \(E=\text{L}^2(X)_\text{ds}\) and call
\((X_\text{kro},T_\text{kro})\) the \textbf{Kronecker subsystem} of
\((X,T)\).

We refer to \((X_\text{kro},\Topology_\text{kro})\) as the
\emph{Kronecker factor} of the system \((X,\Topology)\).

\end{definition}

Kra and Host (\citeproc{ref-kra2007uniformity}{2007}, Proposition 6.1)
and Host and Kra (\citeproc{ref-Host2018NilpotentSI}{2018}, Proposition
24.3) were the sources given for Host
(\citeproc{ref-host2019short}{2019}, Proposition 5). We will compare
these original results with the result we intend to prove.

\begin{proposition}[{Kra and Host, 2007, Proposition
6.1}]\protect\hypertarget{prp-KroneckerGenericSource1}{}\label{prp-KroneckerGenericSource1}

Let \((X,T)\) be a topological dynamical system, \(x_0\in X\) a
transitive point, and \(\Measure\) an invariant ergodic measure on
\(X\). Let \((Y,S)\) be a distal topological dynamical system, \(\nu\)
an invariant measure on \(Y\), and
\(\KroneckerFactorMap:\ (X,\Measure,T)\mapsto(Y,\nu,S)\) a measure
theoretic factor map.

Then there exists a point \(y_0\in Y\) and a sequence of intervals
\(\boldsymbol{I}=(I_j:j\geq1)\) whose lengths tend to infinity such that
for every continuous function \(f\) on \(X\) and every continuous
function \(g\) on \(Y\),
\[\int f(x)(g\circ\KroneckerFactorMap)(x)\ \mathrm{d}\Measure(x) = \lim_{j\rightarrow+\infty}\frac{1}{I_j}\sum_{n\in \CountingMeasure{I_j}}f(T^nx_0)g(s^ny_0). \]

\end{proposition}

Comparing the assumptions for
Proposition~\ref{prp-KroneckerGenericSource0} with
Proposition~\ref{prp-KroneckerGenericSource1}, we have:

{\def\LTcaptype{none} % do not increment counter
\begin{longtable}[]{@{}
  >{\raggedright\arraybackslash}p{(\linewidth - 4\tabcolsep) * \real{0.0577}}
  >{\raggedright\arraybackslash}p{(\linewidth - 4\tabcolsep) * \real{0.4038}}
  >{\raggedright\arraybackslash}p{(\linewidth - 4\tabcolsep) * \real{0.5385}}@{}}
\toprule\noalign{}
\begin{minipage}[b]{\linewidth}\raggedright
\end{minipage} & \begin{minipage}[b]{\linewidth}\raggedright
Proposition~\ref{prp-KroneckerGenericSource0}
\end{minipage} & \begin{minipage}[b]{\linewidth}\raggedright
Proposition~\ref{prp-KroneckerGenericSource1}
\end{minipage} \\
\midrule\noalign{}
\endhead
\bottomrule\noalign{}
\endlastfoot
1 & & \(x_0\in X\) has been specified as a transitive point. \\
2 & \(\Measure\) has been specified as being supported on the closed
orbit of \(x_0\) under the action of \(T\). & \\
3 & \((Y,S)\) is defined as the Kronecker factor \(\KroneckerFactor\).
& \\
\end{longtable}
}

Additionally:

\begin{itemize}
\tightlist
\item
  \(\tilde\Measure\) appears to be a joining between \(\Measure\) and
  \(\nu\).
\item
  \(\boldsymbol{I}\) is the Følner sequence \(\tilde\Phi\).
\end{itemize}

Left to check definitions for:

\begin{itemize}
\tightlist
\item
  (\citeproc{ref-kra2007uniformity}{Kra and Host, 2007}, Theorem 3.1):
  Let \(k\geq1\) be an integer. For a \(k\)-step nilsystem,
  \((X=G/\Gamma,T)\) with Haar measure \(\mu\).

  \begin{itemize}
  \tightlist
  \item
    \(G\) is a \(k\)-step nilpotent Lie group

    \begin{itemize}
    \tightlist
    \item
      \(G\) is a group where, for \(g,h\in G\), we write
      \([g,h]=ghg^{-1}h^{-1}\) for the commutator of \(g\) and \(h\) and
      we write \([A,B]\) for the subgroup spanned by
      \(\{[a,b]:a\in A,b\in B \}\).
    \item
      The commutator subgroups \(G_j\), \(j\geq1\), are defined
      inductively by setting \(G_1=G\) and \(G_{j+1}=[G_j,G]\).
    \item
      Let \(k\geq 1\) be an integer. We say that \(G\) is a \(k\)-step
      nilpotent if \$G\_\{k+1\}is the trivial subgroup.
    \end{itemize}
  \item
    \(\Gamma\) is a discrete cocompact subgroup of \(G\)
  \item
    The compact manifold \(X=G/\Gamma\) is called a \(k\)-step
    nilmanifold

    \begin{itemize}
    \tightlist
    \item
      The group \(G\) acts on \(X\) by left translations and we write
      this action as \((g,x)\mapsto g.x\).
    \end{itemize}
  \item
    \((X,T)\) is \emph{transitive} (admits a dense orbit).

    \begin{itemize}
    \tightlist
    \item
      This is equivalent to being minimal, uniquely ergodic, or
      \((X,\Measure,T)\) being ergodic.
    \end{itemize}
  \end{itemize}
\item
  {[}distal{]}

  \begin{itemize}
  \tightlist
  \item
    previously seen def

    \begin{itemize}
    \tightlist
    \item
      A topological system \((X,T)\) is \textbf{distal} if
      \[\inf\{d(T^nx,T^ny):n\in\mathbb{N}\}>0 \] for all
      \(x\neq y\in X\), where \(d\) is the metric on \(X\).
      (\citeproc{ref-kra2022infinite}{Kra et al., 2022})
    \end{itemize}
  \item
    a def exists for any semigroup action by assuming there's a metric
    compact space
  \item
    a def exists for equicontinuous
  \item
    Discuss why has distal systems been used?

    \begin{itemize}
    \tightlist
    \item
      Is it required for the application of the result?
    \item
      Is it required for the proof?
    \end{itemize}
  \item
    check the Kronecker system in my case is distal
  \end{itemize}
\end{itemize}

\begin{proof}[{Kra and Host, 2007, p. 257}]
{[}Section 4.1 (p.~24) - association of \((X,T)\) and \(x_0\) to a
sequence (the start is just FCP){]}

Section 4.1:

\begin{itemize}
\tightlist
\item
  All FCP?

  \begin{itemize}
  \tightlist
  \item
    Define: separable subalgebra of \(\ell^\infty(\mathbb{Z})\)

    \begin{itemize}
    \tightlist
    \item
      Separable subalgebra of \(\ell^\infty(\mathbb{Z})\) is meant as
      unitary subalgebra of \(\ell^\infty(\mathbb{Z})\) as well as (*)

      \begin{itemize}
      \tightlist
      \item
        ISem28: An amenable group, \(\AmenableGroup\), has a group
        homomorphism \(U:\AmenableGroup\rightarrow\mathscr{U}(H)\) where
        \(\AmenableGroupElement\mapsto U_\AmenableGroupElement\) is a
        unitary representation of \(\Gamma\) on \(H\). The set of all
        unitary operators, \(\mathscr{U}(H)\), is the subset of the set
        of linear isometries \(U:H\rightarrow H\) (such that
        \(||Uf||=||f||\) for all \(f\in H\)) that are surjective (and
        thus bijective).
      \item
        Would the amenable case of the separable subalgebra be the case
        where \(f=U_\gamma\circ r\) where \(r\in\ell^\infty(X)\)?
      \item
        Kra and Host (\citeproc{ref-kra2007uniformity}{2007}) is using a
        functional analysis approach to FCP {[}Attempt following the
        steps with the FCP we covered{]}.
      \item
        For \(\ell^\infty(\AmenableGroup)\), \(a\) could be
        \((\boldsymbol{1}_\AmenableGroupElement:\AmenableGroupElement\in \AmenableGroup)\).
      \item
        Easier to think of \(\ell^\infty(\mathbb{Z})\) as all the
        bounded functions on \(\mathbb{Z}\) and has the supremum norm.
        For \(B\subset\mathbb{Z}\) you can define a function
        \(f\in\ell^\infty(\mathbb{Z})\), i.e., \(\boldsymbol{1}_B\).
      \end{itemize}
    \item
      (*):

      \begin{itemize}
      \tightlist
      \item
        invariant under the shift
      \item
        invariant under complex conjugation
      \item
        closed in \(\ell^\infty(\mathbb{Z})\)
      \item
        {[}separable{]} for the uniform norm \(||\cdot||_\infty\)
      \end{itemize}
    \item
      We mostly consider the case of the separable subalgebra
      \(\mathcal{A}(\boldsymbol{a})\) spanned by a bounded sequence
      \(\boldsymbol{a}=(a_n:n\in\mathbb{Z})\).
    \item
      The shift on \(\ell^\infty(\mathbb{Z})\) is written as \(\sigma\)
      and thus, for a sequence \(\boldsymbol{a}=(a_n:n\in\mathbb{Z})\),
      \(\sigma\boldsymbol{a}=(a_{n+1}:n\in\mathbb{Z})\)
    \item
      We use \(\bar{\boldsymbol{a}}\) to denote the conjugate sequence
      \((\bar{\boldsymbol{a}}:n\in\mathbb{Z})\).
    \end{itemize}
  \end{itemize}
\end{itemize}

{[}If \(Y\) denotes the Kronecker factor of \((X,\Measure,T)\){]}

{[}The sequence of intervals \(I\) plays the same role as the
``Kronecker complete processes'' of {[}BFW{]}{]}

We write \(d_X(\cdot,\cdot)\) and \(d_Y(\cdot,\cdot)\) for distances on
\(X\) and \(Y\), defining the topologies of these spaces.

Subsection 6.1.1 (p.~39)

\begin{itemize}
\tightlist
\item
  \(\mathcal{B}\) closed subalgebra of \(\text{L}^\infty(\Measure)\)
  spanned by \(\mathcal{C}(X)\) and the function \(g\circ\pi\) with
  \(g\in\mathcal{C}(Y)\).

  \begin{itemize}
  \tightlist
  \item
    It is unitary, separable, and invariant under complex conjugation
    and under \(T\).
  \end{itemize}
\item
  \(W\) is the Gelfand spectrum of this algebra.

  \begin{itemize}
  \tightlist
  \item
    Let \(W\) be the Gelfand spectrum of \(\mathcal{B}\) then \(W\)
    consists of the set of unitary homomorphisms from \(\mathcal{B}\) to
    the complex numbers.

    \begin{itemize}
    \tightlist
    \item
      \(W\) will be a ``nice'' space - compact metric
    \item
      Anything we have on \(\mathcal{B}\) will have analogies in \(W\)

      \begin{itemize}
      \tightlist
      \item
        As \(\mathcal{B}\) is invariant under \(T\), then an equivalent
        invariance is in \(W\), see \(R:W\rightarrow W\).
      \item
        Like how \(Tf=f\circ T\) looks at pointwise transformations on
        \(\mathcal{B}\), this goes a step further with \(R\) on \(W\).
      \end{itemize}
    \item
      Can construct the Kronecker factor with the Gelfand spectrum.

      \begin{itemize}
      \tightlist
      \item
        This would require a \(\mathcal{B}\) generated by the
        eigenfunctions, which relates to our previous work on assuming
        an Sigma Algebra of almost-invariant subsets.
      \item
        Compare with ISem28 Exercise 8.5.
      \item
        This method differs to using disintegrations as disintegrations
        start with assuming a sigma algebra of almost-invariant subsets
        whereas Gelfand spectrums starts with all the eigenfunctions
        (which may not completely correspond to almost-invariant
        subsets) and create a sigma algebra from this (thus being more
        general).
      \end{itemize}
    \end{itemize}
  \item
    As \(\mathcal{B}\) is separable, \(W\) is a compact metrizable
    space.
  \item
    By def, there exists an isometric isomorphism of algebras
    \(\Psi:\mathcal{C}(W)\rightarrow\mathcal{B}\).

    \begin{itemize}
    \tightlist
    \item
      For \(\boldsymbol{b}\in\mathcal{B}\), the function
      \(\Psi^{-1}(\boldsymbol{b})\) is called \emph{the function
      associated to} \(\boldsymbol{b}\).
    \item
      The map \(\boldsymbol{b}\mapsto b_0\) is a character of the
      algebra \(\mathcal{B}\). Thus, there exists a point \(w_0\in W\)
      with \$f(w\_0)=\Psi(f)\_0 for all \(f\in\mathcal{C}(W)\)
    \end{itemize}
  \end{itemize}
\item
  The shift on \(\mathcal{B}\) induces a homeomorphism
  \(R:W\rightarrow W\) satisfying \(\Psi(f\circ T)=\Psi(f)\circ R\) for
  all functions \(f\in\mathcal{C}(W)\).

  \begin{itemize}
  \tightlist
  \item
    Therefore, for every \(f\in\mathcal{C}(W)\), \(\Psi(f)\) is the
    sequence \[\Psi(f)=\left(f(T^nw_0):n\in\mathbb{Z} \right). \]
  \item
    If \(f\in\mathcal{C}(W)\) satisfies \(f(T^nw_0)=0\) for all
    \(n\in\mathbb{Z}\), then the sequence given by
    \(\Psi(\boldsymbol{b})=f\) is identically zero and so \(f\) itself
    is identifcally zero.

    \begin{itemize}
    \tightlist
    \item
      It follows that the point \(w_0\) of \(W\) is transitive, meaning
      that its orbit \(\{T^nw_0:n\in\mathbb{Z}\}\) is dense in \(W\).
    \end{itemize}
  \end{itemize}
\item
  The triple \((W,T,w_0)\) is called the \emph{pointed topological
  dynamical system associated to the algebra} \(\mathcal{B}\).
\item
  The inclusion of \(\mathcal{C}(X)\) in \(\mathcal{B}\) induces a
  continuous surjective map \(p:W\rightarrow X\) satisfying
  \(f\circ p=\Psi(f)\) for every continuous function \(f\) on \(X\) and
  we have that \(T\circ p=p\circ R\).

  \begin{itemize}
  \tightlist
  \item
    Similarly, the map \(g\mapsto g\circ\pi\) from \(\mathcal{C}(Y)\) to
    \(\mathcal{B}\) is an isometric homomorphism of algebras and thus
    induces a continuous surjective map \(q:W\rightarrow Y\) satisfying
    \(g\circ q=\Psi(g\circ\pi)\) for all continuous functions \(g\) on
    \(Y\).
  \item
    We have that \(S\circ q=q\circ R\) so \(p:(W,R)\rightarrow(X,T)\)
    and \(q:(W,R)\rightarrow(Y,S)\) are factor maps, in the topological
    sense.
  \end{itemize}
\item
  The map \(f\mapsto\int f\ \mathrm{d}\Measure\) is a positive linear
  form on the algebra \(\mathcal{B}\) and thus, there exists a unique
  probability measure \(\rho\) on \(W\) satisfying
  \[\int f\ \mathrm{d}\Measure=\int\Psi(f)\ \mathrm{d}\rho\qquad\text{for all functions }f\in\mathcal{B}. \]

  \begin{itemize}
  \tightlist
  \item
    As \(\Psi(f\circ T)=\Psi(f)\circ R\) for all \(f\in\mathcal{B}\) and
    \(\mu\) is invariant under \(T\), the measure \(\rho\) is invariant
    under \(R\).
  \item
    As \(\Psi(f)=f\circ p\) for all continuous functions \(f\) on \(X\),
    we have that the image of \(\rho\) under \(p\) is equal to \(\mu\).
  \item
    Therefore, \(p:(W,\rho,R)\rightarrow(X,\mu,T)\) is a measure
    theoretic factor map.

    \begin{itemize}
    \tightlist
    \item
      Moreover, for every function \(f\in\mathcal{B}\),
      \[\int|\Psi(f)|^2\ \mathrm{d}\rho=\int\Psi(|f|^2)\ \mathrm{d}\rho=\int|f|^2\ \mathrm{d}\mu \]
      and the map \(\Psi\) is an isometry from the space \(\mathcal{B}\)
      endowed with the norm \(\text{L}^2(\mu)\) into the space
      \(\text{L}^2(\rho)\).
    \end{itemize}
  \end{itemize}
\item
  Since \(\mathcal{C}(X)\) is dense in \(\mathcal{B}\) under the
  \(\text{L}^2(\mu)\) norm and since \(\Psi(f)=f\circ p\) for
  \(f\in\mathcal{C}(X)\), we have that, for all \(f\in\mathcal{B}\),
  \[\Psi(f)=f\circ p\quad(p\text{-almost everywhere}). \]
\item
  We claim that the map \(p:(W,\rho,R)\rightarrow(X,\mu,T)\) is an
  isomorphism between measure preserving systems.

  \begin{itemize}
  \tightlist
  \item
    Indeed, the range of the map
    \[f\mapsto f\circ p:\text{L}^2(\mu)\rightarrow\text{L}^2(\rho) \] is
    closed in \(\text{L}^2(\rho)\) becuase this map is an isometry, and
    it contains \(\Psi(\mathcal{B})=\mathcal{C}(W)\) and thus it is
    equal to \(\text{L}^2(\rho)\).

    \begin{itemize}
    \tightlist
    \item
      In particular, \((W,\rho,R)\) is ergodic.
    \end{itemize}
  \end{itemize}
\item
  Finally, for every function \(g\in\mathcal{C}(Y)\), we have that
  \(g\circ q=\Psi(g\circ\pi)=g\circ\pi\circ p\) (\(\rho\)-almost
  everywhere) and so \(q=\pi\circ p\) (\(\rho\)-almost everywhere).

  \begin{itemize}
  \tightlist
  \item
    In particular, the image of \(\rho\) under \(q\) is \(\nu\).
  \end{itemize}
\end{itemize}

Subsection 6.1.2 (p.~40)

\begin{itemize}
\tightlist
\item
  Since \(\rho\) is ergocid under \(R\), it admits a generic point
  \(w_1\).

  \begin{itemize}
  \tightlist
  \item
    Recall that this means that, for every \(f\in\mathcal{C}(W)\),
    \[\lim_{j\rightarrow+\infty}\frac{1}{j}\sum_{n=0}^{j-1}f(R^nw_1)=\int f\ \mathrm{d}\rho. \]
  \end{itemize}
\item
  Set \(x_1=p(w_1)\).
\item
  Since \(x_0\) is a transitive point of \(X\), we can choose,
\end{itemize}

\end{proof}

\section{Conclusion}\label{conclusion}

\begin{proposition}[{cf. Host, 2019, Proposition
5}]\protect\hypertarget{prp-KroneckerGeneric}{}\label{prp-KroneckerGeneric}

{Let \((\ShiftSpace,\AmenableGroup)\) be a topological dynamical system
where \(\AmenableGroup\) is an amenable group,
\(\ShiftSpaceElement_0\in\ShiftSpace\), and \(\Measure\) be an ergodic
invariant probability measure supported on the closed orbit of
\(\ShiftSpaceElement_0\) under the action of \(\AmenableGroup\).}

Let \(\KroneckerFactor\) be the Kronecker factor of
\(\FurstenbergSystem\), with factor map
\(\KroneckerFactorMap: \ShiftSpace\rightarrow\KroneckerSpace\).

Let \(\ShiftSpace\times\KroneckerSpace\) be endowed with the group
action of \(\AmenableGroup\times\KroneckerGroup\). Let
\(\tilde\Measure\) be the measure on
\(\ShiftSpace\times\KroneckerSpace\) and image of \(\Measure\) under the
map \(\ShiftSpace\rightarrow\ShiftSpace\times\KroneckerSpace\) where
\(\ShiftSpaceElement\mapsto(\ShiftSpaceElement,\KroneckerFactorMap(\ShiftSpaceElement))\).

Then there exists a Følner sequence \(\tilde{\Folner}\) and a point
\(\KroneckerSpaceElement_0\in\KroneckerSpace\) such that
\((\ShiftSpaceElement_0,\KroneckerSpaceElement_0)\) is generic for
\(\tilde{\Measure}\) along \(\tilde{\Folner}\).

\end{proposition}

\begin{proof}
Let \(d_X(\cdot,\cdot)\) and \(d_Y(\cdot,\cdot)\) for distances on \(X\)
and \(Y\).

Let \(\mathcal{B}\) be the closed subalgebra of
\(\text{L}^\infty(\Measure)\).
\end{proof}

\section{Recent Publications}\label{recent-publications}

\begin{itemize}
\tightlist
\item
  Reflect on amenable result from newer paper within the wider context
  of the paper.
\item
  Write own proof not using Gelfand spectrum.
\end{itemize}

Relatively recently, Charamaras and Mountakis
(\citeproc{ref-charamaras2024findingproductsetsclasses}{2024}) proved
analogous results for amenable groups:

\begin{proposition}[{Charamaras and Mountakis, 2024, Lemma
3.5}]\protect\hypertarget{prp-KroneckerGenericAmenable}{}\label{prp-KroneckerGenericAmenable}

Let \(\AmenableGroup\) be an amenable group, let \(\FurstenbergSystem\)
be an ergodic \(\AmenableGroup\)-system, \(\KroneckerFactor\) be its
Kronecker factor and
\(\KroneckerFactorMap: \FurstenbergSystem\rightarrow\KroneckerFactor\)
be a factor map.

If \(\ShiftSpaceElement\in\ShiftSpace\) is a transitive point, then
there exists a point \(\KroneckerSpaceElement\in\KroneckerSpace\) and a
Følner sequence \(\FurstenbergFolner'\) such that
\[\lim_{N\rightarrow\infty}\frac{1}{\CountingMeasure{\FurstenbergFolner[N]'}}\sum_{\AmenableGroupElement\in\FurstenbergFolner[N]'}f_1(\GroupAction{\AmenableGroupElement}{\ShiftSpaceElement})\cdot f_2(\KroneckerAction{\AmenableGroupElement}{\KroneckerSpaceElement})=\int_\ShiftSpace f_1\cdot(f_2\circ\KroneckerFactorMap)\ \mathrm{d}\mu \]
holds for any \(f_1\in C(\ShiftSpace)\) and
\(f_2\in C(\KroneckerSpace)\).

\end{proposition}

\begin{proposition}[{Charamaras and Mountakis, 2024, Proposition 3.7 (i
\&
(ii))}]\protect\hypertarget{prp-ContinuousKroneckerMap}{}\label{prp-ContinuousKroneckerMap}

Let \(\AmenableGroup\) be an amenable group, let \(\FurstenbergSystem\)
be an ergodic \(\AmenableGroup\)-system, and let
\(\ShiftSpaceElement\in\text{gen}(\Measure,\Folner)\) for some Følner
sequence \(\Folner\). Let \(\KroneckerFactor\) be its Kronecker factor
and
\(\KroneckerFactorMap: \FurstenbergSystem\rightarrow\KroneckerFactor\)
be a factor map, \(\ShiftSpace\times\KroneckerSpace\) be endowed with
the transformation \(T\times S\), and \(\tilde{\Measure}\) be the
joining of \(\FurstenbergSystem\) and \(\KroneckerFactor\).

Then there exists an ergodic extension
\((\ShiftSpace\times\KroneckerSpace, \tilde{\Measure}, T\times S)\) of
\((\ShiftSpace,\Measure,\AmenableGroup)\), a Følner sequence
\(\tilde{\Folner}\), and a point
\((\ShiftSpaceElement,\KroneckerSpaceElement)\in\text{gen}(\tilde{\Measure},\tilde{\Folner})\)
such that:

\begin{enumerate}
\def\labelenumi{\arabic{enumi}.}
\tightlist
\item
  The projection map \(\ProjectionMap[\ShiftSpace]\) is a continuous
  factor map with
  \(\ProjectionMap[\ShiftSpace](\ShiftSpaceElement,\KroneckerSpaceElement)=\ShiftSpaceElement\).
\item
  The projection map \(\ProjectionMap[\KroneckerSpace]\) is a continuous
  factor map where \(\KroneckerFactor\) is the Kronecker factor of
  \((\ShiftSpace\times\KroneckerSpace, \tilde{\Measure}, T\times S)\).
\end{enumerate}

\end{proposition}

\section{Summary}\label{summary}

\begin{theorem}[]\protect\hypertarget{thm-FCP-Kronecker}{}\label{thm-FCP-Kronecker}

Let \(\AmenableGroup\) be an amenable group,
\(A\subset \AmenableGroup\), and \(\Folner\) be a Følner sequence in
\(\AmenableGroup\) such that the limit
\[\delta=\lim_{N\rightarrow\infty}\frac{\CountingMeasure{A\cap\Folner[N]}}{\CountingMeasure{\Folner[N]}}\]
exists.

Then there exists an ergodic system
\((\ShiftSpace,\Measure,\AmenableGroup)\) that is acted on by
\(\AmenableGroup\), a factor map to its ergodic Kronecker factor
\(\KroneckerFactorMap: \FurstenbergSystem\rightarrow\KroneckerFactor\),
Følner sequences \(\FurstenbergFolner\) and \(\tilde{\Folner}\) in
\(\AmenableGroup\), and a point
\((a,\KroneckerSpaceElement_0)\in\ShiftSpace\times\KroneckerSpace\) such
that:

\begin{enumerate}
\def\labelenumi{\arabic{enumi}.}
\tightlist
\item
  We have \(\KroneckerSpaceElement_0=\KroneckerFactorMap(a)\).
\item
  We have clopen sets \(E\subset\ShiftSpace\) and
  \(\tilde{E}\subset\KroneckerSpace\) where
  \(\tilde{E}=\KroneckerFactorMap(E)\) and \begin{align*}
    A&=\{\AmenableGroupElement\in \AmenableGroup:\GroupAction{\AmenableGroupElement}{a}\in E\}
    \\&=\{\AmenableGroupElement\in \AmenableGroup:\KroneckerProductAction{\AmenableGroupElement}{(a,\KroneckerSpaceElement_0)\}\in \tilde{E}}.
  \end{align*}
\item
  \(a\) is generic with \(\Measure\) with respect to
  \(\FurstenbergFolner\) and \(\KroneckerSpaceElement_0\) is generic
  with \(\KroneckerMeasure\) with respect to \(\tilde{Folner}\).
\item
  \(\Measure(E)=\KroneckerMeasure(\tilde{E})\geq\delta\).
\item
  The factor map
  \(\KroneckerFactorMap: \FurstenbergSystem\rightarrow\KroneckerFactor\)
  is continuous.
\end{enumerate}

\end{theorem}

\begin{refremark}
Kronecker factors are \(1\)-step Pronilfactors and the continuous factor
map
\(\KroneckerFactorMap: \FurstenbergSystem\rightarrow\KroneckerFactor\)
shows that a Kronecker factor is a topological pronilfactor.

\label{rem-Kronecker}

\end{refremark}

\section{JdLG-Decomposition}\label{jdlg-decomposition}

\protect\phantomsection\label{refs}
\begin{CSLReferences}{1}{0}
\bibitem[\citeproctext]{ref-charamaras2024findingproductsetsclasses}
Charamaras, D. and Mountakis, A. (2024). \emph{Finding product sets in
some classes of amenable groups}.
https://doi.org/\url{https://doi.org/10.1017/fms.2024.155}.

\bibitem[\citeproctext]{ref-host2019short}
Host, B. (2019). \emph{A short proof of a conjecture of erdös proved by
moreira, richter and robertson}. Available at:
\url{https://arxiv.org/abs/1904.09952}

\bibitem[\citeproctext]{ref-Host2018NilpotentSI}
Host, B. and Kra, B. (2018). `{Nilpotent Structures in Ergodic Theory}'.
\emph{Mathematical Surveys and Monographs}, 236,
\href{https://doi.org/10.1090/surv/236}{https://doi.org/10.1090/surv/236.}

\bibitem[\citeproctext]{ref-ISem28}
Jamneshan, A. and Kreidler, H. (2025). \emph{ISem 28: Ergodic structure
theory and applications}.

\bibitem[\citeproctext]{ref-kra2022infinite}
Kra, B. et al. (2022). \emph{Infinite sumsets in sets with positive
density}. Available at: \url{https://arxiv.org/abs/2206.01786}

\bibitem[\citeproctext]{ref-kra2007uniformity}
Kra, B. and Host, B. (2007). \emph{Uniformity seminorms on
\(\ell^\infty\) and applications}. Available at:
\url{https://arxiv.org/abs/0711.3637}

\end{CSLReferences}




\end{document}
